% \paragraph{Retail and supply-chain decision benchmarks.}
% Recent work on retail and supply-chain decision-making has evolved from isolated subproblems toward integrated, multi-stage environments. Reinforcement-learning benchmarks such as OFCOURSE, GymSC-style simulators, and MARLIM \citep{OFCOURSE, inventory, leluc2023marlimmultiagentreinforcementlearning} model fulfillment and inventory control under demand uncertainty, demonstrating the effectiveness of learned policies in mitigating effects such as demand volatility. However, these environments typically emphasize fixed structures and short- to medium-horizon policies, limiting their ability to evaluate sustained strategic behavior. More recent benchmarks incorporating LLM-based agents, including InvAgent and AIM-Bench \citep{invagentlargelanguagemodel, aimbenchevaluatingdecisionmakingbiases}, begin to explore language-driven reasoning and decision biases, yet still fall short in assessing long-horizon, strategy-aware autonomy in realistic retail settings.

\paragraph{Long-horizon planning benchmarks for LLMs.}
In parallel, a growing body of benchmarks targets long-horizon planning abilities of large language models across structured and interactive environments. PlanBench focuses on classical plan generation, while WebShop, Mind2Web, and ScienceWorld \citep{mind2web, scienceworldagentsmarter5th, 3webshopscalablerealworldweb} evaluate multi-step interaction, error recovery, and adaptive behavior in dynamic settings. More recent benchmarks---such as HeroBench, OdysseyBench, UltraHorizon \citep{utltabench, herobenchbenchmarklonghorizonplanning, odysseybenchevaluatingllmagents}, and explicitly stress extended, interdependent decision sequences that require persistent memory, hierarchical reasoning, and long-term strategy maintenance. Collectively, these benchmarks suggest that while short-horizon tasks are increasingly tractable, robust long-horizon planning and execution remain a central challenge for LLM-based agents.

\paragraph{Long-horizon agent frameworks.}
Recent advances in agent design address these challenges by introducing structured frameworks that decouple high-level planning from low-level execution. Approaches such as Plan-and-Act \citep{erdogan2025planandactimprovingplanningagents} and EAGLET \citep{si2025goalplanjustwish} adopt planner--executor architectures to support hierarchical decision-making, dynamic replanning, and improved execution stability over long horizons. Extensions to multi-agent settings ELHPlan \citep{ling2025elhplanefficientlonghorizontask} and plan-aware context management frameworks PAACE \citep{yuksel2025paaceplanawareautomatedagent} further emphasize task decomposition, explicit strategy representation, and memory-aware context control. Together, these works indicate that scalable long-horizon autonomy increasingly relies on structured planning modules and controlled execution mechanisms rather than monolithic end-to-end policies.